\section{Erd\H{o}s Problem \#624: independent continuation}
\noindent Date: 23 September 2026. Source versions read in full: \texttt{624.tex} in \texttt{gpt\_pro\_5.4}.

\subsection*{1) OBJECTIVE AND SOURCE AUDIT}
The target is $H(n)-\log_2 n\to\infty$, where every Boolean ideal $2^Y$ of rank at least $H(n)$ must receive all $n$ colours. The supplied random-colouring upper bound does not address this lower-bound target. Surjectivity is required, not a rainbow colouring unless $2^{|Y|}=n$.
\subsection*{2) AN INTEGRAL OBSTRUCTION}
Suppose $n=2^h$, $h\ge2$, and rank $h$ suffices. Each ideal $2^Y$ has exactly $n$ members, so its colours must all be distinct. Every two distinct members of
\[
\{\varnothing\}\cup\{\{x\}:x\in X\}
\]
are contained in some such ideal, since their union has size at most two. Hence these $n+1$ sets must have different colours, a contradiction. Thus $H(2^h)\ge h+1$. This is only a weak, already known consequence, not the desired divergent excess; the official discussion records substantially stronger results of Alon.
\subsection*{3) WHY THE NATURAL LINEAR RELAXATION STALLS}
Introduce assignment variables $x_{A,c}\ge0$, with
\[
\sum_c x_{A,c}=1,\qquad \sum_{A\subseteq Y}x_{A,c}\ge1
\quad(|Y|=h,\ c\in X).
\]
An actual colouring is a zero-one solution. But the fractional assignment $x_{A,c}=1/n$ is feasible whenever $2^h\ge n$. Conversely summing the covering inequalities for one $Y$ gives $2^h\ge n$. Therefore this relaxation has exactly the trivial logarithmic threshold. No argument using only linear consequences of these inequalities can prove a stronger lower bound: the explicit fractional witness would contradict it.
\subsection*{4) NEXT STEP AND LIMITATION}
The integral argument uses simultaneous incompatibility across overlapping ideals, a feature absent from that relaxation. A useful next target is a quantitative inequality bounding the number of repeated colours forced inside one ideal when $n$ is just below $2^h$. The argument above forces one repetition only at equality and does not amplify it as $h$ grows.
\subsection*{7) FINAL}
\textbf{UNRESOLVED.} The contribution is a transparent integrality barrier and a self-contained weak lower bound, with no claimed advance over the strongest known bounds.
\noindent Reference checked: \url{https://www.erdosproblems.com/forum/thread/624}.

\subsection*{8) COMPLETION ESTIMATE}
\textbf{COMPLETION: 5\%}
This is a subjective estimate of progress toward the entire stated problem, not a probability that the conjecture or an unverified proof is correct.
